% Part: first-order-logic
% Chapter: syntax-and-semantics
% Section: free-vars-sentences

\documentclass[../../../include/open-logic-section]{subfiles}

\begin{document}

\olfileid{fol}{syn}{fvs}

\olsection{\printtoken{P}{variable} آزاد و \printtoken{P}{sentence}}

\begin{defn}[رخدادهای آزادِ !!a{variable}]
\ollabel{defn:free-occ}
رخدادهای \emph{آزادِ} !!a{variable} در !!a{formula} به‌صورت استقرایی
چنین تعریف می‌شوند:
\begin{enumerate}
\item \indcase*{!A}{$!A$ اتمی است}{همهٔ رخدادهای !!{variable} در
  $\indfrm$ آزادند.}

\tagitem{prvNot}{\indcase{!A}{\lnot !B}{رخدادهای آزادِ !!{variable}
  در $\indfrm$ دقیقاً همان رخدادهای آزاد در $!B$ هستند.}}{}

\item \indcase{!A}{(!B \ast !C)}{رخدادهای آزادِ
  !!{variable} در $\indfrm$ رخدادهای موجود در $!B$
  به‌همراه رخدادهای موجود در~$!C$ هستند.}

\tagitem{prvAll}{\indcase{!A}{\lforall[x][!B]}{رخدادهای آزادِ !!{variable}
  در $\indfrm$ همهٔ رخدادهای آزاد در~$!B$ به‌جز رخدادهای
  $x$ هستند.}}{}

\tagitem{prvEx}{\indcase{!A}{\lexists[x][!B]}{رخدادهای آزادِ !!{variable}
  در $\indfrm$ همهٔ رخدادهای آزاد در~$!B$ به‌جز رخدادهای
  $x$ هستند.}}{}
\end{enumerate}
\end{defn}

\begin{defn}[متغیرهای مقید]
یک رخداد از !!a{variable} در فرمول~$!A$ \emph{مقید} است هرگاه
آزاد نباشد.
\end{defn}

\begin{prob}
در امتداد \olref[fol][syn][fvs]{defn:free-occ}، تعریفی استقرایی از
رخدادهای متغیر مقید ارائه کنید.
\end{prob}

\begin{defn}[دامنه]
\iftag{prvAll}{اگر $\lforall[x][!B]$ رخدادی از یک زیرفرمول در
  فرمول~$!A$ باشد، آنگاه رخداد متناظرِ~$!B$ در~$!A$
  \emph{دامنهٔ} رخداد متناظرِ~$\lforall[x]$ نامیده می‌شود.
  \iftag{prvEx}{برای $\lexists[x]$ نیز به همین ترتیب است.}{}}{اگر
  $\lexists[x][!B]$ رخدادی از یک زیرفرمول در
  فرمول~$!A$ باشد، آنگاه رخداد متناظرِ~$!B$ در~$!A$
  \emph{دامنهٔ} رخداد متناظرِ~$\lexists[x]$ نامیده می‌شود.}

اگر $!B$ دامنهٔ یک رخداد سور
\iftag{prvAll}{$\lforall[x]$\iftag{prvEx}{ یا
    $\lexists[x]$}{}}{$\lexists[x]$} در~$!A$ باشد، آنگاه رخدادهای آزادِ
$x$ در~$!B$ به‌وسیلهٔ همان رخداد سور، در رخداد متناظرِ
\iftag{prvAll}{$\lforall[x][!B]$\iftag{prvEx}{ یا
$\lexists[x][!B]$}{}}{$\lexists[x][!B]$} مقید می‌شوند. می‌گوییم این
رخدادها به‌وسیلهٔ رخداد سور یادشده \emph{مقید شده‌اند}.
\end{defn}

\begin{ex}
فرمول زیر را در نظر بگیرید:
\[
\lexists[\Obj v_0][\underbrace{\Atom{\Obj A^2_0}{\Obj v_0,\Obj v_1}}_{!B}]
\]
$!B$ دامنهٔ $\lexists[\Obj v_0]$ را نمایش می‌دهد.
سور، رخداد $\Obj v_0$ را در $!B$ مقید می‌کند، اما
رخداد $\Obj v_1$ را مقید نمی‌کند. پس در این مورد، $\Obj v_1$
متغیری آزاد است.


اکنون می‌توانیم ببینیم که این سازوکار در یک
!!{formula} پیچیده‌تر، یعنی~$!A$، چگونه عمل می‌کند:
\[
\lforall[\Obj v_0][\underbrace{(\Atom{\Obj A^1_0}{\Obj v_0} \lif
    \Atom{\Obj A^2_0}{\Obj v_0, \Obj v_1})}_{!B}] \lif \lexists[\Obj
  v_1][\underbrace{(\Atom{\Obj A^2_1}{\Obj v_0, \Obj v_1} \lor \lforall[\Obj v_0][\overbrace{\lnot \Atom{\Obj A^1_1}{\Obj v_0}}^{!D}])}_{!C}]
\]
$!B$ دامنهٔ نخستین $\lforall[\Obj v_0]$ است، $!C$ دامنهٔ
$\lexists[\Obj v_1]$ است و $!D$ دامنهٔ دومین
$\lforall[\Obj v_0]$ است. نخستین $\lforall[\Obj v_0]$
رخدادهای $\Obj v_0$ را در~$!B$ مقید می‌کند، $\lexists[\Obj v_1]$
رخداد $\Obj v_1$ را در $!C$ مقید می‌کند و دومین
$\lforall[\Obj v_0]$ رخداد $\Obj v_0$ را در~$!D$ مقید می‌کند.
نخستین رخداد $\Obj v_1$ و چهارمین رخداد $\Obj v_0$ در~$!A$
آزادند. آخرین رخداد $\Obj v_0$ در $!D$ آزاد، اما در $!C$ و~$!A$
مقید است.
\end{ex}

\begin{defn}[جمله]
!!^a{formula}~$!A$ \article{sentence} \emph{!!{sentence}} است اگر و تنها اگر
هیچ رخداد آزاد از !!{variable}s در آن نباشد.
\end{defn}

% add examples!

\end{document}
